\chapter{المجدول وتنقيل السياق}
\label{chap:sched}

يتولى \indextext{المجدول (Scheduler)} توزيع وقت معالجة المعالج بين عدة عمليات جارية بالتناوب،
مما يمنح المستخدم انطباعاً بتنفيذ كافة البرامج متزامنة في نفس الوقت.

\section{تنقيل السياق (Context Switch)}

يُمثل \indextext{تنقيل السياق (Context Switch)} العملية الأساسية لتبديل تنفيذ المعالج من عملية إلى عملية أخرى.
يتطلب تنقيل السياق حفظ السجلات التنفيذية للعملية الحالية بداخل هيكل البيانات المخصص \texttt{struct context}،
ثم تحميل السجلات التنفيذية للعملية الجديدة أو المجدول.

تنفذ النواة تنقيل السياق عبر الدالة المجمّعة \texttt{swtch(old_context, new_context)}.
تحفظ \texttt{swtch} سجلات المكلّف (\texttt{ra}, \texttt{sp}, \texttt{s0}-\texttt{s11}) بداخل بنية السياق القديم،
وتحمل السجلات المقابلة بداخل بنية السياق الجديد، ثم تعود عبر تعليمية \texttt{ret} ليبدأ التنفيذ في المسار الجديد.

تحتفظ كل نواة معالج بمسار تنفيذ خاص بالمجدول المسمى \texttt{scheduler()}.
تنفذ الدالة \texttt{scheduler()} حلقة تكرارية مستمرة تتصفح جدول العمليات بحثاً عن عملية بوضع الجاهزية \texttt{RUNNABLE}،
فتستدعي \texttt{swtch} لنقل التحكم إليها؛ وعندما تنتهي فترة العمل المخصصة للعملية أو تنتظر حَدَثاً، تستدعي العملية الدالة \texttt{yield()} أو \texttt{sleep()} لتعود بالتحكم مجدداً لمسار المجدول.
